% Preamble
\documentclass[../main.tex]{subfiles}

% Document
\begin{document}
\chapter{Methodology}\label{ch:methodology}

As previously mentioned, the practical part of this work will consist of applying 
multi-armed bandits to an e-commerce platform. This section will detail the procedure 
followed to carry out this use case.  

As is common in reinforcement learning applications, instead of directly using a 
dataset to train the implemented models, a simulator or bot was built against which 
the algorithms are used; that is, this bot generates, at each time step (or turn), 
a reward for the action chosen by the algorithm being used. Once the reward is returned, 
the policy of the multi-armed bandit algorithm is updated with respect to its preferences 
over which action to choose in each case, and the process moves to the next turn, where 
the simulator generates a new user, the algorithm suggests the most suitable action for that 
user, and the simulator returns the reward for said action.  

In this case, the user corresponds to a user who has accessed the website, the actions are 
the possible products to recommend, and the reward is whether or not the user clicks on the 
recommended product in order to make a purchase.  

The simulator used is self-developed, and its operation is as follows:  

\begin{itemize}
    \item First, a function is used to generate a set of users who will be the potential 
    customers. These users will have an associated id, age, gender, and preferred product 
    category.
    \begin{itemize}
        \item The id is a number that will allow users to be uniquely identified.
        \item The age will be generated as a random integer between 15 and 70.
        \item he gender will be a random choice between male (M) and female (F).
        \item And the preferred category will be randomly selected from among five 
        categories: technology, fashion, books, sports, and home.
    \end{itemize}
    \item A different function is used to generate the products available on the 
    e-commerce platform. These products will have an associated id, a category, 
    a base click rate received by the product on the platform, and an age factor.
    \begin{itemize}
        \item In this case, the id will also be an integer number that will allow 
        products to be uniquely identified.
        \item The category is the section to which the product belongs. As mentioned 
        for the user generator, the possible categories will be: technology, fashion, 
        books, sports, and home. The simulator is designed such that all categories 
        have the same number of products.
        \item The base click rate defines the base probability (prior to having any 
        type of information about the customer or the time of access to the website) 
        that a given product is clicked if it is recommended to a customer. These 
        probabilities are generated from a uniform distribution with lower bound 0.15 
        and upper bound 0.3, $U\left(0.15, 0.3\right)$.
        \item The age factor estimates how attractive a product is to customers 
        according to their ages. This factor adds a term that follows a linear function 
        to the base click rate (more information about this will be given later). The 
        values of this factor come from random samples of a normal distribution 
        $N\left(0.1, 0.05\right)$, bounded below by 0.02, meaning that if any of the 
        sampled values is below 0.02, it is set to that value, and bounded above by 0.4.
    \end{itemize}
    \item The third and final component of the simulator is the interaction simulator, 
    which is responsible for deciding, for each user-product interaction, whether the 
    user clicked on the recommended product or not.
\end{itemize}

This decision of whether the user clicked on the recommended product is made by sampling 
from a Bernoulli distribution, $Bernoulli\left(p\right)$, whose probability is computed 
as will be explained below:  

For each product, as mentioned previously, there is a base click rate whose values come 
from samples of a continuous uniform distribution $U\left(0.15, 0.3\right)$.  

This base click rate is multiplied by a factor depending on whether the customer’s preferred 
product category matches the product category. This factor is 1.5 if they match and 0.75 if 
they do not.  

It is also multiplied by a gender factor depending on the category of the suggested product. 
If the user is male and the suggested product belongs to the sports or technology categories, 
it is multiplied by a factor of 1.1; if it belongs to the other three categories, the factor 
is 0.9. If the user is female, the factors are reversed: 0.9 for sports and technology products, 
and 1.1 for the other three categories.  

To this click rate, two additional terms are also added, one of which follows a linear 
function depending on the user’s age and the other follows another linear function depending on 
the time at which the user accessed the website.  

Therefore, the final click rate is given by the expression:  

\begin{equation*}
    ctr = base \cdot cat \cdot gen + age + hour
\end{equation*}

where $base$ is the base click rate, $cat \in \left\lbrace 0.75 ,1.5\right\rbrace$ is the 
category-matching factor described above, $gen \in \left\lbrace 0.9 ,1.1\right\rbrace$ is the 
gender factor also previously described, $age$ is the linear term depending on the user’s age 
and $hour$ is the linear term depending on the time at which the user accessed the page.  

The age term is given by the expression:  

\begin{equation*}
    age= \dfrac{a \cdot x}{70}
\end{equation*}

where $a \in N\left(0.1, 0.05\right)$ is the coefficient of the linear regression and $x$ is the 
user’s age. The value of $a$ is also constrained to lie within the interval $\left[0.02, 0.4\right]$.  

The hour term is given by the expression:  

\begin{equation*}
    hour= \dfrac{0.05 \cdot x}{24}
\end{equation*}

where in this case $x$ is the hour of access to the website.  

After designing this simulator, the next step was to use it as a data generator to train the 
reinforcement learning algorithms presented in the state of the art section.  

For the case of the multi-armed bandit algorithms that do not use context ($\epsilon$-Greedy, UCB, and TS), 
this was ignored, whereas for the contextual multi-armed bandit algorithms (LinUCB and LinTS), 
the context was indeed used. However, since this context contains categorical variables (user gender, 
user preferred product category, and product category), these variables had to be transformed, as they 
cannot be directly used to train the linear models underlying these algorithms.  

These variables were encoded using One-Hot encoding, which consists of creating a binary variable for 
each category of one of these categorical variables, such that the variable takes the value 1 if the 
selected data point belongs to that category and 0 otherwise.  

There are other alternative ways of encoding these categorical variables, such as ordinal encoding, in 
which each category of a variable is assigned a different number. However, One-Hot encoding was chosen 
because it does not introduce any bias between the different categories of a variable, unlike ordinal 
encoding, which imposes an ordering on the categories.  

Before performing a comparison of the different models, the hyperparameters of the $\epsilon$-Greedy, Upper 
Confidence Bound, Linear Upper Confidence Bound, and Linear Thompson Sampling algorithms were individually 
optimized in order to obtain the best possible version of each of these models for a fair comparison.  

Note that the hyperparameters of the Thompson Sampling algorithm were not optimized. This is because this 
algorithm lacks numerical hyperparameters similar to those of the other algorithms. In this case, the only 
hyperparameter is the choice of prior distribution used, which was fixed to a beta distribution.  

After this hyperparameter optimization, all models were trained simultaneously, such that the user selected 
at each time step of the problem was the same for all algorithms. For this user, whether they would click on 
each of the different products available on the platform was simulated, assuming each of them was recommended 
to them. Once this was done, the models were called to produce their recommendations, and the result of the 
suggested product was returned. Using the outcome of each algorithm’s recommendation, they were updated 
(each one only with the result of its own suggested product), and the process moved on to the next iteration.  

To evaluate the performance of all these algorithms, both cumulative reward and cumulative regret were used. 
In addition to these metrics, accuracy (number of times the best possible action was suggested divided by the 
number of time steps) and the proportion of visits that resulted in the user clicking on the suggested product 
(success rate) were also analyzed.  
\end{document}